\subsection{An added derivative has its own starting value}
\label{sec:added-instrument-origins}

Increasing the instrument volume introduces a second measurement origin.
IFRS~9.B6.5.19 retains the existing instrument's history and measures the added
instrument's change from the date of rebalancing. If the company pays for an
instrument with a positive fair value, that payment establishes an asset under
IFRS~9.5.1.1. Only changes after designation contribute to its hedge gain
\citep[IFRS~9, paragraphs 5.1.1 and B6.5.19]{eu20231803}.

Return to the original 100-unit exposure designated at price 80, measured at
price 90 with derivative gain 1200, reserve 1000 and cumulative profit 200.
Here the item volume stays at 100. Suppose changed circumstances require adding
an eligible external derivative A, acquired and designated together for its
fair value of 200. The risk-management objective is unchanged, the revised
ratio is eligible and the documentation is updated. Debit derivative A 200
and credit cash 200. Total derivative carrying value is now 1400, but the
cumulative hedge gain remains 1200: the company acquired the additional 200
of value with cash. Reserve and profit remain 1000 and 200.

At price 95, suppose the original instrument is worth 1400 and A is worth 220.
The original instrument has gained 1400 since its zero-value designation;
A has gained $220-200=20$. The combined gain is therefore 1420, while total
carrying value is 1620. The item loss is $100(95-80)=1500$, giving reserve
1420 and cumulative profit zero. Since the preceding measurement, derivative
value rose by 220, reserve rose by 420 and earnings fell by 200. The entry
reverses the earlier ineffectiveness. Substituting carrying value 1620 for
gain 1420 would give reserve 1500 and profit 120, both incorrect for this
designation history.

The general relation separates the acquired value from later change. Let
$b_j$ be instrument $j$'s fair value when it is acquired and designated, and
$V_j(t)$ its current value. The original zero-value instrument has $b_0=0$.
With no settlement or removal, define total acquired value $B$, cumulative
hedge gain $D_t$ and total derivative carrying value $V_t$ by
\[
 B=\sum_j b_j,\qquad D_t=\sum_j\bigl(V_j(t)-b_j\bigr),
 \qquad V_t=B+D_t.
\]
An acquisition increases both $B$ and $V_t$ by the same amount, leaving $D_t$
unchanged. Later changes in value alter $D_t$. The reserve calculation uses
$D_t$ and the cumulative hedged-item change $H_t$:
\[
 R_t=\operatorname{sgn}(D_t)\min(|D_t|,|H_t|),\qquad I_t=D_t-R_t,
\]
within the same offsetting, supplied-valuation conditions as the preceding
example. The item layers keep their original designation prices; an instrument
purchase does not reset them.

Acquire a second additional derivative for 50 after that measurement. At the next
observation, price 97, let the original instrument and the two additions have values
1900, 260 and 60. Carrying value is 2220; cumulative gain is
$1900+(260-200)+(60-50)=1970$. The item loss is 1700, so reserve is 1700
and cumulative profit 270. With opening cash 20000, the two acquisitions leave
cash 19750. Cash plus derivatives is $19750+2220=21970$, equal to capital
20000, reserve 1700 and profit 270. The acquisition payments belong in that
reconciliation even though they created no earnings.

Different instruments can later have opposite carrying signs. If, at price
79, the values become $-200$, 100 and 20, the company has derivative assets
120 and a derivative liability 200. Their signed sum is $-80$, whereas
cumulative hedge loss is $-80-250=-330$. The item gain is 100, leaving a
reserve loss 100 and cumulative loss 230. Cash 19750 plus assets 120 equals
liability 200 plus equity 19670. Keeping the separate positions preserves
their meaning; the signed total alone gives no permission to offset assets
and liabilities in the statements.

The documentary \code{AddedInstrumentLayers/1} report checks each instrument's
identity, acquired value, designation date and subsequent marks. Missing marks,
unfunded purchases and unsupported designation changes are refused. Marks and
eligibility are supplied findings, with all effects of the whole instruments
included. This acquisition example assumes equal transaction price and fair
value, no fees and no earlier local book position. An already-owned derivative,
excluded valuation components, instrument removals and settlement need their
own histories and methods. The unchanged objective and updated designation
remain substantive conditions; changing quantities merely to suppress observed
ineffectiveness is not justified by these calculations
\citep[IFRS~9, paragraphs 6.2.1--6.2.6, B6.5.12--15 and B6.5.21]{eu20231803}.
